% magnetism-body.tex
% Shared content for both track versions of the magnetism equation sheet.
% \input this from a file that has already set up the preamble (packages,
% colors, \eqtab/\R/\band macros) and opened \begin{document} with its own
% title block.
%
% Everything here is common to PHYS 230 (algebra) and PHYS 235 (calc). The
% guiding idea of this sheet: every magnetic law is an electric law with one
% extra ingredient inserted -- a cross product with a velocity (forces/fields
% of moving charges) or with a current element (sources). Section 1 states
% that pattern explicitly; the rest of the sheet is built to keep showing it.
%
% Track-specific material (the Biot--Savart differential / integral) lives
% only in magnetism-calc.tex, after this \input.
%
% Ampere's law is NOT on either sheet (removed 2026-11-12): PHYS 230 never
% covered it and Seth is not testing it in PHYS 235 either. Both tracks work
% from the given result formulas (wire, loop, solenoid).
%
% Cross-references use \label/\ref, never a typed section number.

%========================================================
\section{The pattern: electric $\to$ magnetic \quad {\normalfont\small (how to read the rest of this sheet)}}\label{sec:pattern}
\eqtab{
  \band{insert a $\vec v\times$}
  \R{Force on a charge\newline\footnotesize(same $q$; the field now has to be crossed with $\vec v$, not just multiplied)}{\vec F=q\vec E \quad\longrightarrow\quad \vec F=q\vec v\times\vec B}
  \R{Field of a point charge (Coulomb's law) $\to$ field of a \emph{moving} point charge\newline\footnotesize(same $1/r^2$; direction is now set by the right-hand rule on $\vec v\times\hat r$, not just radially out)}{\vec E=\dfrac{1}{4\pi\varepsilon_0}\dfrac{q}{r^2}\hat r \quad\longrightarrow\quad \vec B=\dfrac{\mu_0}{4\pi}\dfrac{q\vec v\times\hat r}{r^2}}
  \band{the source swap: $q\vec v \to I\,\Delta\vec L$}
  \R{A single moving charge generalizes to a continuous current\newline\footnotesize(this is exactly how the moving-point-charge law above becomes the Biot--Savart law, Sec.~\ref{sec:sources} -- a current element is just a lot of moving charge packed into a short wire segment)}{q\vec v \quad\longrightarrow\quad I\,\Delta\vec L}
  \band{electric dipole $\to$ magnetic dipole}
  \R{Dipole moment, torque, and orientation energy all carry straight across\newline\footnotesize($\vec p\to\vec m$, $\vec E\to\vec B$ -- see Sec.~\ref{sec:dipoles} for the full set side by side)}{\vec p \to \vec m\,,\qquad \vec E \to \vec B}
}

%========================================================
\section{Magnetic force on charges \& currents \quad {\normalfont\small (Knight 29.1--29.3, 29.7--29.8)}}\label{sec:force}
\eqtab{
  \band{fundamental}
  \R{Force on a moving charge\newline\footnotesize($\perp$ to both $\vec v$ and $\vec B$; use the right-hand rule; zero if $\vec v \parallel \vec B$; compare $\vec F=q\vec E$, Sec.~\ref{sec:pattern})}{\vec F = q\,\vec v \times \vec B\,,\qquad |\vec F| = |q|\,vB\sin\theta}
  \R{Force on a straight current-carrying wire, length $L$}{\vec F = I\,\vec L \times \vec B\,,\qquad |\vec F| = BIL\sin\theta}
  \band{derived}
  \R{Charge in a uniform field ($\vec v \perp \vec B$)\newline\footnotesize(circular motion; magnetic force does no work, so $v$ is constant; here $m$ is the particle's mass)}{r = \dfrac{mv}{|q|B}}
  \R{Period / cyclotron frequency\newline\footnotesize(independent of $v$ and $r$)}{T = \dfrac{2\pi m}{|q|B}\,,\qquad f = \dfrac{|q|B}{2\pi m}}
  \R{Velocity selector (crossed $\vec E$, $\vec B$)}{v = \dfrac{E}{B}}
}

%========================================================
\section{Sources of magnetic field \quad {\normalfont\small (Knight 29.4--29.5)}}\label{sec:sources}
\eqtab{
  \band{fundamental}
  \R{Field of a single moving point charge\newline\footnotesize(the magnetic twin of Coulomb's law, Sec.~\ref{sec:pattern} -- same $1/r^2$, direction by the right-hand rule on $\vec v\times\hat r$)}{\vec B = \dfrac{\mu_0}{4\pi}\dfrac{q\vec v \times \hat r}{r^{2}}}
  \R{Chop a current into pieces, add up\newline\footnotesize(each short segment $\Delta\vec L$ of current $I$ acts like a swarm of moving charges, $I\Delta\vec L$ standing in for $q\vec v$ -- see Sec.~\ref{sec:pattern})}{\Delta\vec B = \dfrac{\mu_0}{4\pi}\,\dfrac{I\,\Delta\vec L \times \hat r}{r^{2}}\,,\qquad \vec B = \textstyle\sum_i \Delta\vec B_i}
  \band{derived --- given results, use these directly}
  \R{Long straight wire (distance $r$)\newline\footnotesize(field circles the wire; $B \propto 1/r$)}{B = \dfrac{\mu_0 I}{2\pi r}}
  \R{Center of a circular loop (radius $R$, $N$ turns)}{B = \dfrac{\mu_0 N I}{2R}}
  \R{Inside a long solenoid\newline\footnotesize(uniform; independent of radius; $\approx 0$ outside)}{B = \mu_0 n I}
}

%========================================================
\section{Magnetic dipoles \& torque \quad {\normalfont\small (Knight 29.5, 29.9)}}\label{sec:dipoles}
\eqtab{
  \band{fundamental}
  \R{Dipole moment of a current loop\newline\footnotesize(compare the electric dipole moment $\vec p=q\vec d$)}{\vec m = N I \vec A \qquad (m = NIA)}
  \R{Torque on a dipole in a field\newline\footnotesize(twists $\vec m$ toward $\vec B$; max at $\theta = 90^\circ$, zero when aligned; compare $\vec\tau=\vec p\times\vec E$)}{\vec\tau = \vec m \times \vec B\,,\qquad |\vec\tau| = mB \sin\theta}
  \band{derived}
  \R{Orientation energy\newline\footnotesize(minimum when $\vec m \parallel \vec B$; compare $U=-\vec p\cdot\vec E$)}{U = -\vec m \cdot \vec B = -mB \cos\theta}
  \R{Net force on a dipole in a \emph{uniform} field}{\vec F_{\text{net}} = 0 \quad (\text{torque only})}
}

%========================================================
\section{Magnetic flux \& electromagnetic induction \quad {\normalfont\small (Knight 30.1--30.7)}}\label{sec:induction}
\eqtab{
  \band{fundamental}
  \R{Magnetic flux through a loop\newline\footnotesize(uniform field: just area times the component of $\vec B$ along the normal)}{\Phi_B = BA\cos\theta}
  \R{Faraday's law --- induced emf\newline\footnotesize(only a \emph{changing} flux induces an emf)}{\mathcal{E} = -N\,\dfrac{\Delta\Phi_B}{\Delta t}}
  \R{Lenz's law}{\text{induced current opposes the change in flux}}
  \band{derived}
  \R{Motional emf (bar of length $L$, speed $v$, $\vec v \perp \vec B$)}{\mathcal{E} = BLv}
  \R{Induced current and dissipated power}{I = \dfrac{\mathcal{E}}{R}\,,\qquad P = \dfrac{\mathcal{E}^{2}}{R} = F_{\text{pull}}\,v}
  \R{Induced electric field (inside a region of changing $B$, radius $r$ from center)\newline\footnotesize(circulates even with no wire present; direction by Lenz's law, as if a loop were there)}{E(r) = \dfrac{r}{2}\left|\dfrac{dB}{dt}\right|}
}
